\documentclass{article}

\usepackage{amsmath, amsthm, amssymb, amsfonts}
\usepackage{thmtools}
\usepackage{graphicx}
\usepackage{setspace}
\usepackage{geometry}
\usepackage{float}
\usepackage{hyperref}
\usepackage[utf8]{inputenc}
\usepackage[english]{babel}
\usepackage{framed}
\usepackage[dvipsnames]{xcolor}
\usepackage{tcolorbox}

\colorlet{LightGray}{White!90!Periwinkle}
\colorlet{LightOrange}{Orange!15}
\colorlet{LightGreen}{Green!15}

\newcommand{\HRule}[1]{\rule{\linewidth}{#1}}

\declaretheoremstyle[name=Theorem,]{thmsty}
\declaretheorem[style=thmsty,numberwithin=section]{theorem}
\tcolorboxenvironment{theorem}{colback=LightGray}

\declaretheoremstyle[name=Proposition,]{prosty}
\declaretheorem[style=prosty,numberlike=theorem]{proposition}
\tcolorboxenvironment{proposition}{colback=LightOrange}

\declaretheoremstyle[name=Principle,]{prcpsty}
\declaretheorem[style=prcpsty,numberlike=theorem]{principle}
\tcolorboxenvironment{principle}{colback=LightGreen}

\setstretch{1.2}
\geometry{
    textheight=9in,
    textwidth=5.5in,
    top=1in,
    headheight=12pt,
    headsep=25pt,
    footskip=30pt
}

% ------------------------------------------------------------------------------

\begin{document}

% ------------------------------------------------------------------------------
% Cover Page and ToC
% ------------------------------------------------------------------------------

\title{ \normalsize \textsc{}
		\\ [2.0cm]
		\HRule{1.5pt} \\
		\LARGE \textbf{\uppercase{Relatório do trabalho\\ Loja de Carros}
		\HRule{2.0pt} \\ [0.6cm] \LARGE{Ferramentas da Internet} \vspace*{10\baselineskip}}
		}
\date{}
\author{\textbf{Autores:} \\
		Tiago 1\\
		Kaique de Freitas 2\\
		Otavio Leone 3\\
        Lucas Estevam 4\\
        Mateus Augusto 5\\
        MARIANA - 2026}

\maketitle
\newpage

% ------------------------------------------------------------------------------

\section{Introdução}
O presente \textbf{trabalho} apresenta um \textit{Minimundo de um banco de dados} para uma loja de carros onde será possível gerenciar veículos, clientes, funcionários, vendas, fornecedorers, compras, financiamentos, \textit{test-drives}, manutenções e estoque.

\section{Minimundo} 
Uma loja de carros necessita de um sistema de banco de dados para gerenciar informações relacionadas a veículos, clientes, funcionários, vendas, fornecedores, compras, financiamentos, test-drives, estoque e manutenção. A loja oferece uma variedade de veículos e cada tipo é gerenciado por um funcionário. Os funcionários fazem as vendas dos veículos. Os clientes compram os veículos. Caso o cliente compre um veículo, ele pode realizar um financiamento como forma de pagamento. Cada cliente pode realizar um test-drive nos veículos e realizar manutenções em seus carros que possuem. Caso o veículo seja comprado, é preciso renovar o estoque dele. Para isso o funcionário entra em contato com o fornecedor para comprar um novo veículo para o estoque da loja. No sistema, a entidade Pessoa se relaciona com Cliente e Funcionário por meio de uma generalização, onde cada pessoa pode ser especializada como cliente ou funcionário. Um funcionário gerencia vários veículos, enquanto cada veículo é gerenciado por apenas um funcionário. Da mesma forma, um funcionário realiza várias vendas, mas cada venda é realizada por apenas um funcionário. Vários clientes compram carros ao longo do tempo, porém cada venda está associada a apenas um cliente. Cada venda está associada a um único veículo, e cada veículo pode ser vendido apenas uma vez. Nem toda venda possui financiamento, mas cada financiamento está vinculado a uma única venda. Um cliente pode realizar vários test-drives, e um veículo pode ser utilizado em diversos test-drives. No caso das manutenções, um cliente pode realizar várias manutenções em seus veículos, e um veículo pode passar por várias manutenções ao longo do tempo. Um fornecedor pode estar associado a várias compras realizadas pela loja, enquanto cada compra é feita com apenas um fornecedor. Cada compra inclui um ou mais veículos, mas cada veículo está associado a uma única compra. Por fim, cada veículo está associado a um único registro no estoque, e uma compra repõe um ou mais itens no estoque. No projeto proposto, cada Pessoa possui CPF, nome, idade, e-mail e telefone. O Cliente, que é uma especialização de Pessoa, possui um ID de cliente, data de nascimento e histórico de compras. Já o Funcionário, também especialização de Pessoa, possui um ID de funcionário, cargo e salário. Cada Veículo tem um ID de veículo, marca, modelo, ano e preço. O Fornecedor possui um ID de fornecedor, nome, telefone, e-mail e endereço. Cada Compra apresenta um ID de compra, data, valor total, forma de pagamento e status. Cada Venda possui um ID de venda, data, valor total, forma pagamento e status. O Financiamento contém um ID de financiamento, valor financiado, número de parcelas, taxa de juros e status. Cada Test-drive possui um ID de test-drive, data, hora, observações e status. A Manutenção é composta por ID de manutenção, descrição, data, custo e tipo de serviço. Por fim, o Estoque possui um ID de estoque, quantidade, valor, capacidade máxima e status.



\newpage


\end{document}
